\begin{textsample}{POS Dim 2 – global\_north\_2025\_02 – Score 19.00 – global\_north\_2025\_02/121132286.txt}  \label{ex:f2_pos_439}
Israel ' s \textbf{relentless} \textbf{siege} stifles Gaza ' s recovery despite ceasefire Ahmed Al-Majayda tries to rebuild his humble home with clay and scraps of material ( Mohamed Solaimane ).
Ahmed Al-Majayda, a 30-year-old father from Khan Younis, Gaza, is rebuilding his home with clay.
The cement he had once hoped to use is now a luxury, priced at an astronomical \$250 for a 50-kilogram bag, compared to just \$7 before Israel ' s devastating assault.
After 14 months of displacement, Al-Majayda is piecing together a small room for his wife and three \textbf{children} in the Al-Mawasi \textbf{neighbourhood}, using scraps salvaged from the rubble of his destroyed home."Not only did I have to remove the rubble by hand, but I also waited for building materials to enter Gaza.
But they never came,"Al-Majayda said, his face marked with exhaustion.
Across Gaza, the blockade ' s cruelty is felt in every corner.
Families like Al-Majayda ' s struggle with the scarcity of basic construction materials, while \textbf{hospitals} operate @ @ @ @ @ @ @ @ @ @, and the \textbf{health} sector is left gasping for life."This is not the life I imagined when I heard about the ceasefire,"Al-Majayda said."The occupation has kept us waiting.
All of our lives depend on what they allow us to receive."\textbf{Hospitals} hit the hardest The fragile ceasefire, while bringing temporary relief, has failed to lead to meaningful improvements in Gaza.
As more than half a million displaced people began returning to \textbf{northern} Gaza last month, they were met with shock at the extent of the destruction.
Despite promises of aid and rebuilding materials, Israel has systematically restricted the entry of vital \textbf{supplies} since the deal took effect on January 19. \textbf{Hospitals} and clinics are suffering under the weight of insufficient resources."The \textbf{health} sector is in a catastrophic state and has not seen any change since the ceasefire came into effect, as the necessary needs to restore life to this sector were not \textbf{available},"Dr Marwan al-Hams, director of field @ @ @ @ @ @ @ @ @ @ allowed 74 \textbf{trucks} into Gaza -- a cruel faction of what is needed."It ' s a joke.
The essential \textbf{medical} \textbf{supplies} that have trickled in amounted to less than 0.03 percent of the sector ' s needs,"he stated.
Reuters A \textbf{truck} moves from the Gaza Strip through the Rafah border crossing on the Egyptian side, amid a ceasefire between Israel and Hamas, in Rafah, Egypt, February 11, 2025 ( Reuters/Amr Abdallah Dalsh ).
With \textbf{hospitals} operating on skeleton crews and \textbf{medical} \textbf{supplies} scarce, doctors face overwhelming pressure."We have only been able to \textbf{reopen} one clinic in Rafah and part of the \textbf{Indonesian} \textbf{Hospital} in the north.
But the bulk of our \textbf{hospitals} remain closed, and our needs continue to be ignored,"Dr Al-Hams said.
The result is a dire situation for the sector. \textbf{Patients} continue to flood in with urgent needs, from \textbf{surgery} to chronic care -- many of their afflictions and illnesses directly tied to Israel ' s \textbf{bombardment} of the strip. @ @ @ @ @ @ @ @ @ @, diabetes, and chronic stress among those seeking treatment.
According to Al-Hams,"the \textbf{shortage} in Gaza ' s \textbf{health} sector remains as much as 80 percent of its needs."He also criticised Israel for failing to uphold its commitment to \textbf{medical} evacuations."The agreement called for 150 \textbf{patients} to leave Gaza daily, but fewer than 50 have been allowed to travel each day."Reuters Some \textbf{children}, like this Palestinian girl, are amongst those able to receive much needed \textbf{medical} care during the ceasefire ( Reuters/Amr Abdallah Dalsh ).
With Israel reducing the number of aid \textbf{trucks}, hindering reconstruction efforts, and restricting \textbf{patient} travel, the temporary halt in fighting has done little to alleviate the ongoing suffering.
For Al-Hams, the immediate focus is securing generators and \textbf{fuel} to ensure \textbf{hospitals} can continue to operate.
Israel ' s political game Israel ' s manipulation of \textbf{humanitarian} aid serves as a ruthless political tool.
Talal Abu Rakba, a political analyst, describes Israel ' s obstruction of aid as @ @ @ @ @ @ @ @ @ @ absolute victory"-- not through military might alone, but by crushing Palestinian resilience."This is not just a delay, but a form of negotiating pressure,"he explains, comparing the agreed-upon terms with the reality of aid entering Gaza.
According to the Gaza Government Media Office, only 8,500 truckloads of \textbf{supplies} have been allowed into Gaza -- well below the 12,000 that were promised since the ceasefire began.
Abu Rakba suggests that Israel aims to extend the first phase of the ceasefire while avoiding deeper political and military concessions."Israel ' s approach is to use these tools to force Hamas into submission,"he says, noting that while Hamas remains militarily weakened, it still holds significant influence due to its ideological roots.
The analyst believes Israel will likely continue restricting aid and hindering reconstruction efforts as a way to exert further pressure on Hamas."This is part of Israel ' s broader strategy to create conditions that push Palestinians into voluntary displacement, despite Egypt ' s opposition to @ @ @ @ @ @ @ @ @ @ struggle to rebuild their lives amid ruins, the aid that trickles through Israeli checkpoints is designed for survival, not recovery."Israel has made its agenda clear : they are forcing Hamas and the people of Gaza into a corner,"Abu Rakba concluded.
The ceasefire is little more than a pause in active \textbf{bombardment}.
The \textbf{siege} continues -- silent, calculated, and devastating.

% matched lemmas: available, bombardment, child, fuel, health, hospital, humanitarian, indonesian, medical, neighbourhood, northern, patient, relentless, reopen, shortage, siege, supply, surgery, truck
\end{textsample}
